\documentclass{article}
\usepackage{url,amsmath,setspace,amssymb,fullpage,color,hyperref}
\usepackage[
n,
operators,
advantage,
sets,
adversary,
landau,
probability,
notions,
logic,
ff,
mm,
primitives,
events,
complexity,
asymptotics,
keys, lambda]{cryptocode}

\newcommand{\heading}[5]{
   \renewcommand{\thepage}{#1-\arabic{page}}
   \noindent
   \begin{center}
   \framebox[\textwidth]{
     \begin{minipage}{0.9\textwidth} \onehalfspacing
       {\bf CS 601.641/441 -- Blockchains and Cryptocurrencies} \hfill #2

       {\centering \Large #5

       }\medskip

       {\it #3 \hfill #4}
     \end{minipage}
   }
   \end{center}
}

\newcommand{\scribe}[4]{\heading{#1}{#2}{}{ #4}{ #3}}

\begin{document}

\heading{3}{Spring 2026}{Instructor: Matthew Green}{TA: Dorian Liu}{Assignment 2: Ethereum --- Phase 3 Submission}

\section*{Student Information}

\begin{tabular}{l l}
\textbf{Name (string):} & \texttt{TODO: Firstname Lastname} \\[6pt]
\textbf{JHED ID:}        & \texttt{TODO: abc123} \\
\end{tabular}

\section*{On-Chain Information}

\begin{tabular}{l p{10cm}}
\textbf{Wallet Address:}        & \texttt{TODO: 0x...} \\[6pt]
\textbf{Attacker Contract:}     & \texttt{TODO: 0x...} \\[6pt]
\textbf{Attack Tx Hash:}        & \texttt{TODO: 0x...} \\
\end{tabular}

\vspace{1em}

\noindent\textbf{Note:} The TA will verify your submission by:
\begin{enumerate}
    \item Checking that \texttt{dao.grades(yourAddress)} returns 100 on Sepolia
    \item Checking that \texttt{dao.ta()} returns your address on Sepolia
    \item Verifying the attack transaction executed successfully in a single transaction
    \item Running the autograder against your \texttt{Attacker.sol} source code
\end{enumerate}

\vspace{1em}

\section*{Short Answer}

\begin{enumerate}

    %% ============================================================
    %% Q1: Aave Lending Protocol
    %% ============================================================
    \item \textbf{Q1: Aave Lending Protocol.}

    Aave is one of the largest decentralized lending protocols in DeFi. Read about Aave's architecture and answer the following.

    \begin{itemize}
        \item[(a)] In Aave, lenders deposit assets and receive \textbf{aTokens} (e.g., depositing USDC gives you aUSDC). Explain how aTokens represent a lender's position. How does a lender earn interest --- does someone call a function to update their balance, or does it happen automatically? Explain the mechanism.
        \item[(b)] Aave supports both \textbf{overcollateralized borrowing} and \textbf{flash loans}. Compare the two: who provides the safety guarantee in each case, and what happens when the guarantee is violated? (i.e., what happens when a borrower's collateral drops below the threshold vs.\ when a flash loan is not repaid?)
        \item[(c)] Aave uses a \textbf{health factor} to determine when a position should be liquidated. If a user deposits 10 ETH (at \$2,000/ETH) with a liquidation threshold of 80\%, and borrows \$12,000 USDC:
        \begin{enumerate}
            \item What is their health factor? (Formula: $H = \frac{\text{collateral} \times \text{liquidation threshold}}{\text{debt}}$)
            \item At what ETH price does the position become liquidatable ($H < 1$)?
            \item Why do liquidators receive a bonus (typically 5--10\%) for liquidating unhealthy positions? What would happen without this incentive?
        \end{enumerate}
        \item[(d)] Aave V3 introduced \textbf{Efficiency Mode (E-Mode)} which allows higher loan-to-value ratios for correlated assets (e.g., borrowing USDC against DAI at 97\% LTV instead of the usual 75\%). Why is it safer to allow higher leverage for correlated assets? What risk remains even for stablecoin-to-stablecoin borrowing?
    \end{itemize}

    \paragraph{Solution.}

    \vspace{4em}

    %% ============================================================
    %% Q2: Terra/LUNA Collapse
    %% ============================================================
    \item \textbf{Q2: The Terra/LUNA Collapse.}

    In May 2022, the Terra ecosystem collapsed, wiping out approximately \$40 billion in value. Read about the Terra/LUNA collapse and answer the following.

    \begin{itemize}
        \item[(a)] Explain how TerraUSD (UST) maintained its \$1 peg under normal conditions. What was the role of LUNA in the mint/burn mechanism?
        \item[(b)] Describe the \textbf{death spiral} that occurred in May 2022. What triggered the initial depeg, and how did the mint/burn mechanism make it worse instead of restoring the peg?
        \item[(c)] Compare UST's algorithmic stability mechanism to a \textbf{collateralized stablecoin} like DAI (MakerDAO). What fundamental design difference explains why DAI survived market crashes that killed UST?
    \end{itemize}

    \paragraph{Solution.}

    \vspace{4em}

    %% ============================================================
    %% Q3: Flash Loans
    %% ============================================================
    \item \textbf{Q3: Flash Loans.}

    \begin{itemize}
        \item[(a)] Explain how a flash loan works at a high level. What makes it different from a traditional loan? Why is no collateral required?
        \item[(b)] A flash loan borrower calls \texttt{flashLoan(1000000)}. Walk through what happens step-by-step inside the \texttt{LendingPool} contract, including the callback mechanism. What happens if the borrower does not repay?
        \item[(c)] Name two legitimate (non-exploit) use cases for flash loans in DeFi.
    \end{itemize}

    \paragraph{Solution.}

    \vspace{4em}

    %% ============================================================
    %% Q4: Your Attack
    %% ============================================================
    \item \textbf{Q4: Your Attack Sequence.}

    Describe the exploit you implemented in your \texttt{Attacker.sol} contract.

    \begin{itemize}
        \item[(a)] What is the vulnerability in the \texttt{GovernanceDAO} contract? Which function is buggy, and why? How would a secure implementation differ?
        \item[(b)] List the exact sequence of contract calls your \texttt{Attacker} makes during the \texttt{onFlashLoan} callback, in order. For each call, explain why it is necessary and what state it changes.
        \item[(c)] Why must the entire exploit happen within a single transaction? What would go wrong if the steps were split across multiple transactions?
    \end{itemize}

    \paragraph{Solution.}

    \vspace{4em}

\end{enumerate}

\end{document}
